\chapter{Circular Motion and Continuous Change}
\label{ch:ch05-circular-motion}

\section{When Angles Move}

The previous chapter extended trigonometry beyond right triangles by
defining sine and cosine on the unit circle: an angle determined a point on
the circle, and the coordinates of that point supplied the values of the
trigonometric functions. Everything studied there was nonetheless static ---
an angle was specified, a point was identified, a value was computed --- and
many of the phenomena that motivate trigonometry in the first place are not
static at all. A wheel turns, a planet moves through its orbit, a gear
rotates, and the hands of a clock sweep continuously through time; in each
case the angle itself is changing, and it is more natural to treat it as a
function of time, $\theta = \theta(t)$, than as a fixed quantity handed down
in advance. The central question of this chapter is simple: what
mathematics emerges once a point is allowed to move continuously around a
circle? The answer will connect geometry, motion, and the trigonometric
functions into a single framework.

\section{Arc Length}

Suppose a point travels along a circle of radius $r$, tracing a portion of
the circumference called an \emph{arc}. The distance traveled along the
circle is the \emph{arc length}, denoted $s$, and the relationship between
$s$ and the central angle $\theta$ depends on the unit in which $\theta$ is
measured. When radians are used, that relationship turns out to be
remarkably simple.

Recall from Chapter~4 that one radian is defined as the angle subtending an
arc whose length equals the radius, so that $\theta = s/r$ is not a theorem
but the very definition of radian measure. Solving for $s$ gives $s =
r\theta$ immediately, and the simplicity of this relationship is exactly
why radians, rather than degrees, play the central structural role in the
mathematics that follows: degrees are convenient for measurement, but
radians are convenient for structure.

\begin{theorem}
For a circle of radius $r$, the arc length subtended by a central angle
$\theta$, measured in radians, is $s = r\theta$.
\end{theorem}

\begin{proof}
By the definition of radian measure, $\theta = s/r$. Multiplying both sides
by $r$ gives $s = r\theta$.
\end{proof}

\section{Sector Area}

A \emph{sector} is the portion of a circle enclosed by two radii and the arc
connecting them, and its area should evidently increase as either the
radius or the angle increases. The formula is obtained by comparing a
sector to the full circle: a complete revolution sweeps out $2\pi$ radians
and encloses area $\pi r^2$, so a sector of angle $\theta$ should occupy the
same fraction of the full area that $\theta$ occupies of $2\pi$,
\[
\frac{A}{\pi r^2} = \frac{\theta}{2\pi}.
\]

\begin{theorem}
The area of a sector with radius $r$ and central angle $\theta$, measured in
radians, is $A = \dfrac12 r^2 \theta$.
\end{theorem}

\begin{proof}
Multiplying both sides of $A/(\pi r^2) = \theta/(2\pi)$ by $\pi r^2$ gives
$A = (\theta/2\pi)\pi r^2$, which simplifies to $A = \tfrac12 r^2\theta$.
\end{proof}

\begin{example}
A sprinkler rotates through an angle of $2\pi/3$ radians and sprays water
to a distance of $6$ meters. Find the arc length swept by the edge of the
spray and the area watered. Using $s = r\theta$,
\[
s = 6\left(\frac{2\pi}{3}\right) = 4\pi \text{ meters,}
\]
and using $A = \tfrac12 r^2\theta$,
\[
A = \frac12 (6^2)\left(\frac{2\pi}{3}\right) = 12\pi \text{ square meters.}
\]
\end{example}

\section{Angular Velocity}

Introducing time allows the rate at which an angle changes to be measured
directly.

\begin{definition}
The \emph{angular velocity} of a rotating object, denoted $\omega$, is the
rate at which its angle changes. For uniform rotation, $\omega = \theta/t$.
\end{definition}

Angular velocity is commonly expressed in radians per second, radians per
minute, or revolutions per minute, and it measures how rapidly an angle
changes regardless of the size of the object doing the rotating: a small
gear and a large gear turning at the same angular velocity complete the
same fraction of a revolution in the same time, even though a point on the
larger gear travels much farther.

\section{Linear Velocity}

That last observation motivates a second, related quantity. A point on the
rim of a bicycle wheel travels farther than a point near the hub even
though both complete the same rotation in the same time, so angular
velocity alone does not capture the actual distance traveled. Starting from
$s = r\theta$ and dividing both sides by time $t$ gives
\[
\frac{s}{t} = r\cdot\frac{\theta}{t},
\]
and recognizing $s/t$ as the \emph{linear velocity} $v$ and $\theta/t$ as
the angular velocity $\omega$ yields the relationship between them.

\begin{theorem}
For uniform circular motion, $v = r\omega$. A larger radius produces a
larger linear velocity even when the angular velocity remains unchanged.
\end{theorem}

\begin{algorithmproc}{How to Relate Angular and Linear Velocity}
Given a point rotating about a center, identify the radius $r$ from the
center to the point and the angular velocity $\omega$ of the rotation.
Apply $v = r\omega$ to obtain the linear velocity, and check that the units
of $r$ and $\omega$ are consistent (typically meters and radians per
second) before reporting the result.
\end{algorithmproc}

\begin{example}
A carousel rotates at $0.5$ radians per second, and a rider sits $4$ meters
from the center. Using $v = r\omega$,
\[
v = (4)(0.5) = 2 \text{ m/s,}
\]
so the rider travels two meters along the circular path every second.
\end{example}

\begin{practicenow}
A wheel rotates at $3$ radians per second and has radius $0.4$ meters.
Find the linear velocity of a point on the rim.
\end{practicenow}

\section{Uniform Circular Motion}

This section arrives at the central idea of the chapter. Consider a point
moving around the unit circle with constant angular velocity $\omega$,
beginning at angle $0$. After time $t$ the angle has become $\theta =
\omega t$, and from Chapter~4 the point on the unit circle at angle
$\theta$ has coordinates $(\cos\theta, \sin\theta)$. Substituting $\theta =
\omega t$ gives the position at time $t$ as $(\cos\omega t, \sin\omega t)$:
the horizontal coordinate of the moving point changes according to a cosine
function of time, and the vertical coordinate changes according to a sine
function of time. Sine and cosine are, in other words, the projections of
uniform circular motion onto two perpendicular axes, and this observation
is among the most important in the entire subject, since it is what will
eventually let an oscillation --- a swinging pendulum, a vibrating string, a
sound wave --- be understood as the shadow of some circle turning at
constant speed.

\section{Coordinates of Uniform Circular Motion}

The same idea extends immediately to motion beginning at an arbitrary
starting angle $\theta_0$ rather than at $0$. After time $t$, the angle has
become $\theta = \omega t + \theta_0$.

\begin{theorem}
A point moving uniformly around the unit circle with angular velocity
$\omega$ and initial angle $\theta_0$ has coordinates
\[
x(t) = \cos(\omega t + \theta_0), \qquad y(t) = \sin(\omega t + \theta_0).
\]
\end{theorem}

\begin{proof}
From Chapter~4, a point at angle $\theta$ has coordinates $(\cos\theta,
\sin\theta)$. Uniform rotation with initial angle $\theta_0$ gives $\theta =
\omega t + \theta_0$, and substituting this into the coordinate formulas
gives $x(t) = \cos(\omega t + \theta_0)$ and $y(t) = \sin(\omega t +
\theta_0)$.
\end{proof}

\section{Applications}

Circular motion appears throughout science and engineering, from rotating
machinery, gears, and flywheels to planetary orbits, electrical generators,
sound waves, and signal processing. The formulas of this chapter reveal
that every periodic oscillation can, in principle, be understood as the
shadow of some circular motion projected onto an axis, an idea that will
eventually lead to Fourier analysis in Chapter~19, where complicated
signals are decomposed into collections of such circular motions added
together. For now, the chapter offers a new interpretation of sine and
cosine themselves: no longer merely ratios read off a triangle, but
descriptions of motion unfolding in time.

\section{Looking Ahead}

This chapter introduced motion into trigonometry, developing formulas for
arc length and sector area, relating angular velocity to linear velocity,
and discovering that sine and cosine arise naturally as the coordinates of
uniform circular motion. The functions themselves, however, remain in a
sense still hidden: their role in describing motion around a circle has
been established, but their shape, plotted on ordinary coordinate axes, has
not yet been examined directly. What happens when $y = \sin\theta$ or $y =
\cos\theta$ is plotted as a curve, with $\theta$ running along the
horizontal axis instead of wrapping around a circle? The next chapter takes
up the graphs of the trigonometric functions, and shows how the periodic
motion studied here appears once it is unrolled and viewed through time.

\section{Exercises}
\begin{exercise}
A circle has radius $9$ centimeters. Find the arc length and sector area
corresponding to a central angle of $5\pi/6$ radians.
\end{exercise}

\begin{exercise}
A Ferris wheel has radius $15$ meters and completes one revolution every
$40$ seconds. Find its angular velocity in radians per second, and the
linear velocity of a rider at the rim.
\end{exercise}

\begin{exercise}
A point moves uniformly around the unit circle with angular velocity
$\omega = \pi/3$ radians per second, starting at $\theta_0 = \pi/4$. Write
its position $(x(t), y(t))$ as a function of $t$, and find its position at
$t = 2$ seconds.
\end{exercise}

\begin{exercise}
Two gears of different radii are connected so that their rims move at the
same linear velocity. Explain, using $v = r\omega$, why the gear with the
smaller radius must have the larger angular velocity.
\end{exercise}
